\chapter{Introdução}

A avaliação de modelos é um dos pilares centrais do Aprendizado de Máquina, orientando a seleção de modelos, a interpretação de resultados e o avanço metodlógico da área. No contexto supervisionado, métricas de desempenho podem ser diretamente calculadas a partir de rótulos verdadeiros, uma vez que são conhecidos, permitindo comparações mais objetivas entre desempenho de modelos. Entretanto, no aprendizado não supervisionado, particularmente em tarefas de agrupamento, a ausência de rótulos reais torna o problema de avaliação significativamente mais desafiador. Nesses cenários, a qualidade de uma partição deve ser inferida a partir da própria estrutura dos dados ou da comparação indireta entre modelos, levantando questões sobre robustez, interpretabilidade e generalidade das medidas utilizadas.

Tradicionalmente, a validação de agrupamentos é realizada por meio é realizada por meio de medidas internas ou externas. Medidas externas como o Rand Index ajustado dependem da existência de uma partição de referência e são amplamente utilizadas em experimentos controlados. Contudo, tais medidas tornam-se inviáveis em aplicações reais, onde o objetivo do agrupamento desconhece o rótulo verdadeiros. Medidas internas, por sua vez, avaliam propriedades como coesão e separação dos clusters, mas são sensíveis à métrica de distância adotada, podendo apresentar vieses em relação a determinadas estruturas e frequentemente produzem ranqueamentos contraditórios entre modelos. Além disso, ambas as abordagens negligenciam uma dimensão fundamental: a variabilidade de dificuldade entre instâncias, tratando todos os pontos de dados com uma dificuldade constante.

Nesse contexto, esta dissertação adota a Teoria de Resposta ao Item (TRI) como arcabouço teórico para modelar avaliação no cenário de Aprendizado de Máquina. Originalmente desenvolvida na psicometria para estimar habilidades latentes de respondentes e dificuldades de itens, a TRI permite analisiar desempenho considerando explicitamente a interação entre habilidade e dificuldade. Em anos recentes, modelos baseados em TRI foram adaptados para avaliar sistemas de Inteligência Artificial em tarefas supervisionadas, possibilitando estimar habilidade de modelos e identificar instâncias mais desafiadoras ou podendo ser consideradas como ruidosas. Essa perspectiva oferece uma alternativa às métricas tradicionais, ao introduzir uma modelagem probablística que incorpora um disciminante e uma dificuldade associadas as instâncias, fornecendo interpretações mais refinadas sobre o desempenho.

A primeira contribuição desta dissertação consiste no aprimoramento do uso para o modelo $\beta^{4}$-IRT, uma extensão do $\beta^{3}$-IRT voltada para respostas contínuas limitadas entre 0 e 1. O modelo trata explicitamente o problema de não-identificabilidade associado ao parâmetro de discriminação, fatorando-o em componentes de sinal e magnitude e adotando uma estratégia de estimação que melhora a recuperação de parâmetros. Essa reformulação mostra-se particularmente relevante em cenários de Aprendizado de Máquina, nos quais discriminações negativas podem aparecer em razão de instãncias ruidosas ou ambíguas.

A segunda contribuição é o método CLAIRE (CLuster Agreement-based Item REsponses), que reformula a avaliação de agrupamentos como um problema de estimação de habilidade latente via TRI. Em vez de depender do rótulo real da instância, o método constrói uma matriz de resposta baseada na concordância entre modelos quanto ao agrupamento de pares de instâncias. Sob essa perspectiva, modelos mais informativos são aqueles que tendem a concordar entre si nas instâncias mais difíceis, enquanto instãncias difíceis são aquelas nas quais apenas modelos de alta habilidade apresentam um consenso. Experimentos demosntram que essa abordagem produz ranqueamentos altamente correlacionados com medidas externas tradicionais e permanece robusta mesmo quando partições aleatórias são adicionadas ao conjunto de modelos avaliados.

Ao integrar avanços metodológicos na modelagem de discriminação com uma nova fomulação para validação de agrupamentos, esta dissertação estabelece a habilidade latente como princípio unificador para avaliação de métodos não supervisionados. Dessa forma, contribui tanto par ao desenvolvimento teórico de métricas de avaliação quanto para aplicações práticas em cenários onde a ausência de rótulos é a regra, e não a exceção.